% =============================================================================
%  text/THE ARGUMENT/fourfold-strategy.tex
%  Theory subsection: demonstrates a 2x2 typology table with booktabs.
%  理论小节示例：演示 2x2 分类表格（三线表）。
% =============================================================================

Combining the ruler's authority security (secure vs.\ fragile) with the
priority given to state-building (high vs.\ low) yields a fourfold typology of
ruling strategies, summarized in Table~\ref{tab:fourfold}.

% 表 1：四类典型策略的 2x2 分类（booktabs 三线表示例）
\begin{table}[htbp]
  \centering
  \caption{A Fourfold Typology of Ruling Strategies}
  \label{tab:fourfold}
  \begin{tabular}{p{3.2cm} p{4.6cm} p{4.6cm}}
    \toprule
    & \textbf{State-building priority: high} & \textbf{State-building priority: low} \\
    \midrule
    \textbf{Authority secure}    & Foster elite cohesion;
                                   invest in institutions. &
                                   Consolidate personal networks;
                                   status quo politics. \\
    \addlinespace
    \textbf{Authority fragile}   & Co-opt elites; share rents to
                                   buy survival. &
                                   Fragment elite coalitions;
                                   purge potential challengers. \\
    \bottomrule
  \end{tabular}
  \begin{tablenotes}
    \footnotesize
    \item \textit{Note:} This table is a placeholder demonstrating the
    \texttt{booktabs} + \texttt{threeparttable} table format. Replace the
    content with your own typology.
  \end{tablenotes}
\end{table}

We expect the empirical case to display a transition across these cells over
time, as the ruler's authority security varies (Section~\ref{sec:ruler-power-dynamics}).
